\section{训练细节、成本与另类架构：为什么这是系统工程}

访谈中“训练细节和成本”虽然章节较短，但它是理解整期的钥匙。罗福莉不断强调，Agent 时代所有路径都要回到端到端完成率、速度和成本。一个模型在 benchmark 上更强，不代表它在 Agent 框架中更优；一个模型更大，也不代表它在所有环节都值得调用。真正的系统工程，是把模型能力、任务阶段、工具调用、延迟、价格和可靠性放在同一张表里做决策。

\begin{importantbox}{成本不是财务尾项，而是模型能力的一部分}
如果一个 Agent 任务需要几十轮模型调用、工具执行和环境反馈，那么单次调用成本会被长程任务放大。便宜小模型、专用模型、多模型路由、缓存、验证器和框架编排，都会直接影响系统能不能上线。
\end{importantbox}

\subsection{术语消化：训练与成本词表}

\noindent\begin{tabularx}{\textwidth}{p{0.23\textwidth}p{0.31\textwidth}X}
\toprule
术语 & 解决的问题 & 与本期关系 \\
\midrule
端到端完成率 & 整个任务最终是否完成 & Agent 产品不能只看局部动作或单轮回答。 \\
调用路由 & 不同环节选择不同模型 & 用大模型做关键推理，小模型做高频低风险环节。 \\
验证器 & 判断任务结果是否正确 & 是 RL feedback 和可靠性评估的关键。 \\
Rollout 成本 & 生成多条轨迹的算力与 API 成本 & 决定后训练能否 scale。 \\
Latency & 用户等待时间和任务总耗时 & 影响生产力工具是否可用。 \\
Scaffold & Agent 外部框架、工具编排和状态管理 & 能弥补模型短板，也能放大顶尖模型上限。 \\
\bottomrule
\end{tabularx}

\subsection{另类架构为什么重要}

另类架构在这里不是“奇技淫巧”，而是资源约束下的可扩展研究。DeepSeek-V2、MoE、attention 改造等方向，本质上都在问同一个问题：如何在给定算力和硬件条件下获得更好的单位成本智能。Agent 时代这个问题更尖锐，因为长程任务会反复调用模型，任何单步效率差距都会被放大。

\begin{warningbox}{不要把另类架构只当论文创新}
如果一个架构无法 scale 到工业模型，无法部署到真实推理系统，或无法降低端到端任务成本，它对 Agent 产品的价值就有限。访谈强调的“工业级水准”，正是要求架构创新最终进入可运行系统。
\end{warningbox}

\subsection{本章小结}

训练细节、成本和另类架构共同说明：Agent 时代的技术路线不是“更大模型”单变量竞赛，而是端到端系统效率竞赛。模型结构、后训练、框架和推理成本必须一起优化。

